% W4i26_New_Predictions.tex
% DFD 17-Agent Research Programme — Iteration 26
% Agents 1 (Formalism), 2 (Lab/MEMS), 3 (Clocks), 5 (Lensing/GW), 13 (Particle), 14 (Astro), 15 (New Predictions)
% Date: 2026-03-21

\documentclass[11pt,a4paper]{article}
\usepackage[utf8]{inputenc}
\usepackage{amsmath,amssymb,amsfonts}
\usepackage{hyperref}
\usepackage{booktabs}
\usepackage{longtable}
\usepackage{geometry}
\geometry{margin=2.5cm}

\title{DFD i26 New Predictions \& Experimental Frontiers\\Agents 1, 2, 3, 5, 13, 14, 15}
\author{DFD 17-Agent Research Programme}
\date{2026-03-21}

\begin{document}
\maketitle
\tableofcontents
\newpage

%=============================================================================
\section{Agent 1: Formalism --- Disformal Framework Consequences}
%=============================================================================

\subsection{Mandate}
Explore consequences of the A2 disformal upgrade. What new predictions arise from the disformal factor $\Sigma(y) = 1/\sqrt{1+y}$? Map DFD onto DHOST classification. Identify new axiom structure.

\subsection{New Literature}

\begin{enumerate}

\item \textbf{Langlois (2025), arXiv:2412.04135} --- ``DHOST theories as disformal gravity: From black holes to radiative spacetimes.'' Habilitation thesis reviewing full DHOST theory space. Degeneracy classification (Classes Ia, Ib, IIa, IIb, III). Stability under disformal field redefinition (DFR). Exact rotating BH solutions. Radiative GW solutions in DHOST. \textbf{Grade: A.}

DFD's disformal Lagrangian maps to \textbf{Class Ia} DHOST (the ``healthy'' class):
\begin{itemize}
\item Conformal factor: $\Omega(\psi) = \mu(\hat{a}/a_0)$ --- field-dependent.
\item Disformal factor: $\Gamma(\psi, X) = \Sigma(y) = (1+y)^{-1/2}$ where $y = X/a_0^2$, $X = \partial_\mu\psi\,\partial^\mu\psi$.
\item Class Ia guarantees: no Ostrogradski ghost, tensor speed $c_T = c$ (if $\alpha_T = 0$), stable cosmological perturbations.
\end{itemize}

\item \textbf{Takahashi \& Motohashi (2026), arXiv:2601.09164} --- Ghost-free scalar-tensor with third-order derivatives. Extends GDH and U-DHOST. \textbf{Grade: B.} DFD has NO third-order derivatives, placing it in the simpler (safer) subclass.

\item \textbf{Kozak (2026), arXiv:2603.08401} --- Palatini disformal Horndeski. Reduces to kinetic gravity braiding on-shell. \textbf{Grade: B.} Alternative geometric interpretation: DFD's $n = e^\psi$ could be viewed as defining a Palatini-type independent connection $\hat{\Gamma} = \Gamma + C(\psi, \partial\psi)$ where $C$ encodes the refractive deviation.

\item \textbf{Takahashi et al.\ (2025), arXiv:2509.05027} --- GLPV unique GW signatures. The $f^5$ induced spectrum from scalar-scalar-tensor interaction requires GLPV disconnected from Horndeski via disformal transformation. \textbf{Grade: B.} DFD's disformal factor $\Sigma(y)$ could generate such an interaction if the mapping goes beyond standard Horndeski. Investigation needed.

\end{enumerate}

\subsection{New Predictions from Disformal Upgrade}

The disformal factor $\Sigma(y) = 1/\sqrt{1+y}$ introduces:

\begin{enumerate}
\item \textbf{Velocity-dependent screening}: Unlike conformal coupling (depends on $|\nabla\psi|$), disformal coupling depends on $\dot{\psi}^2$ (kinetic energy). Systems with rapid scalar-field evolution experience stronger screening. \textbf{Prediction P30 (NEW)}: In merging galaxy clusters (high $\dot{\psi}$), the effective gravitational coupling deviates from static predictions by $\sim\Sigma(y) - 1 \approx -y/2$ for small $y$.

\item \textbf{Disformal lensing correction}: Photon geodesics in the disformal metric $\tilde{g}_{\mu\nu}$ acquire a velocity-dependent correction to the deflection angle:
\begin{equation}
\delta\alpha_\mathrm{disf} = -\frac{1}{2}\frac{(\partial_t\psi)^2}{a_0^2}\,\alpha_\mathrm{GR}
\end{equation}
For galaxy-cluster lensing where $\partial_t\psi \sim H_0\,\Phi_N/c$, this is $\sim 10^{-12}\,\alpha_\mathrm{GR}$ --- far too small to detect currently.

\item \textbf{Modified gravitational slip}: The ratio $\eta = \Phi/\Psi$ acquires a disformal correction:
\begin{equation}
\eta_\mathrm{DFD} = 1 + \frac{\Sigma'(y)\,y}{1 + \Sigma(y)\,y} \approx 1 - \frac{y}{2(1+y)}
\end{equation}
At galaxy scales ($y \ll 1$): $\eta \approx 1$ (GR-like). In deep MOND regime ($y \sim 1$): $\eta \approx 0.75$. \textbf{Prediction P31 (NEW)}: Euclid measurements of the gravitational slip $\eta$ in the weak-lensing/galaxy-clustering cross-correlation should find $\eta$ consistent with 1 at $z > 1$ but potentially deviating at $z < 0.5$.
\end{enumerate}

\textbf{Grade: A. Disformal upgrade generates two new testable predictions (P30, P31).}

%=============================================================================
\section{Agent 2: Lab/MEMS --- Atom Interferometry and Torsion Balance}
%=============================================================================

\subsection{Mandate}
Track AION, MAGIS-100, ZAIGA, torsion balance experiments. Assess sensitivity to DFD scalar field.

\subsection{New Literature}

\begin{enumerate}

\item \textbf{MAGIS-100 status (2025--2026)} --- Under construction at Fermilab. 100-meter vertical atom interferometer using Sr atoms. Mid-band GW sensitivity: 0.01--3 Hz. UK AION collaboration contributing optics, cameras, DAQ. Funded by Moore Foundation (\$9.8M) + DOE QuantISED + UKRI. \textbf{Grade: B.}

DFD scalar field $\psi$ produces a differential phase shift in atom interferometers:
\begin{equation}
\delta\phi_\psi = k_\mathrm{eff}\,\Delta\psi\,T^2
\end{equation}
where $\Delta\psi$ is the scalar field difference over the 100-m baseline, $T$ the interrogation time. For Earth's gravitational gradient: $\Delta\psi \approx g\,L/c^2 \approx 10^{-14}$. With $T \sim 1$ s, $k_\mathrm{eff} \sim 10^7$ m$^{-1}$: $\delta\phi_\psi \sim 10^{-7}$ rad. Current sensitivity: $\sim 10^{-4}$ rad. \textbf{Three orders of magnitude short.} MAGIS-100 cannot yet detect DFD's scalar field directly.

\item \textbf{Eot-Wash group (2026), arXiv:2602.02815} --- ``An Improved Torsion Balance Test of the Equivalence Principle Towards the Sun.'' Year-long dataset (July 2024--July 2025). Tests long-range interactions between normal matter and dark matter using the Sun as source. \textbf{Grade: B.} DFD predicts no EP violation (the scalar field couples universally via $n = e^\psi$), so this result is a consistency check. Any detection of EP violation would challenge DFD.

\item \textbf{Qin et al.\ (2024), Phys.\ Dark Univ.\ 44, 101481} --- ``Constraining light scalar field with torsion-balance gravity experiments.'' HUST-18 bounds: photon coupling $\Lambda_\gamma = 7\times10^{17}$ GeV, electron coupling $\Lambda_e = 1\times10^{17}$ GeV for scalar mass $10^{-9}$--$10^{-4}$ eV. \textbf{Grade: B.} DFD's scalar $\psi$ is massless (or ultralight), so these constraints apply at the boundary of the mass range. The coupling constraints are not yet sensitive enough to reach DFD's predicted coupling scale.

\item \textbf{Biercuk et al.\ (2025), arXiv:2512.06333} --- ``Testing the weak equivalence principle for nonclassical matter with torsion balances.'' Proposes using superposition of energy eigenstates as test masses. \textbf{Grade: C.} Conceptually novel but far from implementation.

\end{enumerate}

\subsection{Assessment}
Lab-scale tests remain 2--3 orders of magnitude from DFD sensitivity thresholds. MAGIS-100 (operational $\sim$2027) and AION-10 (Oxford, $\sim$2028) will improve atom interferometric sensitivity by $\sim$100$\times$, potentially bringing DFD's scalar field within reach at the $\sim$1$\sigma$ level.

\textbf{Grade: B (no change from i25). Long-term prospect improving.}

%=============================================================================
\section{Agent 3: Clocks --- Thorium-229 Nuclear Clock Timeline}
%=============================================================================

\subsection{Mandate}
PRIORITY: When exactly will nuclear clock comparisons be available? Which groups? Update K = 5900$\pm$2300 sensitivity.

\subsection{New Literature}

\begin{enumerate}

\item \textbf{Tiedau et al.\ (2026), Nature s41586-025-09999-5; arXiv:2507.01180} --- ``Frequency reproducibility of solid-state Th-229 nuclear clocks.'' Key results:
\begin{itemize}
\item At 195 K: frequency reproducibility of 220 Hz (fractionally $1.1\times10^{-13}$) for two differently doped $^{229}$Th:CaF$_2$ crystals over 7 months.
\item Optimal working temperature: $196\pm5$ K (first-order thermal sensitivity vanishes).
\item With in-situ temperature co-sensing using different quadrupole-split lines, systematic shift reducible below $10^{-18}$ fractional frequency uncertainty.
\end{itemize}
\textbf{Grade: A.} This is a landmark result. The $10^{-13}$ reproducibility is already sufficient for a crude $\alpha$-variation test. The $10^{-18}$ target would make Th-229 competitive with the best optical clocks.

\item \textbf{Banerjee et al.\ (2025), Nature Commun.\ s41467-025-64191-7; arXiv:2407.17300} --- ``Fine-structure constant sensitivity of the Th-229 nuclear clock transition.'' Using state-resolved spectroscopy at $10^{-12}$ precision: $K = 5900 \pm 2300$. Quadrupole moment change: $\Delta Q_0/Q_0 = 1.791(2)\%$. Three orders of magnitude enhancement over atomic clocks. \textbf{Grade: A.}

\item \textbf{Ye group / JILA (2025)} --- VUV frequency comb interrogation of $^{229}$Th:CaF$_2$. Transition frequency: $2{,}020{,}407{,}384{,}335(2)$ kHz. First direct comparison with Sr atomic clock. kHz-level precision. \textbf{Grade: A.}

\item \textbf{Thirolf et al.\ (2025), Nat.\ Sci.\ Rev.\ nwaf083} --- ``The ticking of thorium nuclear optical clocks: a developmental perspective.'' Reviews TU Wien (crystal growth), PTB (VUV laser), JILA (VUV comb). Identifies path to $10^{-19}$ fractional uncertainty with trapped single $^{229}$Th$^{3+}$ ions. \textbf{Grade: B.}

\item \textbf{Ponce et al.\ (2025), Phys.\ Rev.\ Lett.\ 134, 113801} --- ``Temperature Sensitivity of a Thorium-229 Solid-State Nuclear Clock.'' Characterizes thermal shift coefficients. \textbf{Grade: B.}

\end{enumerate}

\subsection{Timeline for Nuclear-Optical Clock Comparison (P29)}

\begin{center}
\begin{tabular}{lll}
\toprule
\textbf{Year} & \textbf{Milestone} & \textbf{Group} \\
\midrule
2024 & First nuclear-atomic comparison (kHz) & JILA (Ye) \\
2025 & Reproducibility $1.1\times10^{-13}$ & TU Wien / PTB \\
2026 & Sub-Hz spectroscopy expected & JILA / PTB \\
2027--2028 & $10^{-16}$ fractional uncertainty (solid-state) & TU Wien / JILA \\
2028--2030 & Single-ion $^{229}$Th$^{3+}$ clock ($10^{-18}$--$10^{-19}$) & PTB / NIST \\
\bottomrule
\end{tabular}
\end{center}

\subsection{DFD Prediction P29 Update}
DFD predicts a scalar-field-dependent shift in the nuclear transition:
\begin{equation}
\frac{\delta\omega_\mathrm{nuc}}{\omega_\mathrm{nuc}} = K\,\frac{\delta\alpha}{\alpha} = K\,\kappa\,\psi
\end{equation}
where $\kappa$ is the DFD coupling of $\alpha$ to $\psi$. With $K = 5900 \pm 2300$ and current clock precision $10^{-13}$:
\begin{itemize}
\item Detectable $\delta\alpha/\alpha > 10^{-13}/5900 \sim 2\times10^{-17}$ --- already competitive.
\item At $10^{-18}$ precision: $\delta\alpha/\alpha > 2\times10^{-22}$ --- would probe DFD's predicted annual modulation from Earth's orbital eccentricity ($\delta\psi \sim 3\times10^{-10}\,e\,\cos\theta$).
\end{itemize}

\textbf{P29 Grade: A (upgraded). Nuclear clocks are on track for DFD-relevant sensitivity by 2027--2028.}

%=============================================================================
\section{Agent 5: Lensing/GW --- GW Propagation in Disformal Gravity}
%=============================================================================

\subsection{Mandate}
GW propagation in disformal gravity. Speed of GW constraints. Lensing predictions.

\subsection{New Literature}

\begin{enumerate}

\item \textbf{de Rham \& Tolley (2025), arXiv:2511.19762} --- (Shared with Agent 7.) GW speed constraints eliminate quartic galileon, constrain disformal coupling $M_\gamma \gtrsim 10$ MeV. \textbf{Grade: B.} DFD safe: $\alpha_T = 0$ by construction.

\item \textbf{Takahashi et al.\ (2025), arXiv:2509.05027} --- (Shared with Agent 1.) GLPV unique $f^5$ induced GW signature. Third derivatives in tensor source. \textbf{Grade: B.}

\item \textbf{Saha \& Modak (2025), arXiv:2504.04331} --- ``Shedding Light on Gravity: Black Hole Shadows and Lensing Signatures in Lorentz Gauge Theory.'' BH shadow deviations in modified gravity. Lensing signatures as observational discriminators. \textbf{Grade: C.} Different theory class but methodology (shadow + lensing joint analysis) applicable to DFD disformal BH solutions.

\item \textbf{Barausse et al.\ (2025), Nature Astronomy s41550-025-02695-4} --- ``The future ability to test theories of gravity with black-hole shadows.'' EHT + ngEHT projections. Shadow size constrains post-Newtonian deviations inaccessible to weak-field tests. \textbf{Grade: B.} DFD's disformal BH solutions (from Langlois 2412.04135) predict specific shadow modifications. The ngEHT (2030s) could test these.

\item \textbf{Bernardo \& Said (2025), Gen.\ Relativ.\ Gravit.\ s10714-025-03437-7} --- ``Black Holes as Laboratories: Tests of General Relativity.'' Comprehensive review of BH-based gravity tests. \textbf{Grade: C.}

\end{enumerate}

\subsection{DFD GW Propagation Summary}
\begin{itemize}
\item Tensor speed: $c_T = c$ exactly ($\alpha_T = 0$). GW170817 satisfied.
\item Scalar breathing mode: predicted with amplitude $h_s/h_+ \sim \psi/\Phi_N \sim \mu(x)/(1+\mu(x))$. In deep MOND ($x \gg 1$): $h_s/h_+ \to 1$ (equal amplitude). In Newtonian ($x \ll 1$): $h_s/h_+ \to x$ (suppressed).
\item BH shadow: disformal correction $\sim y/(1+y)$ where $y = (\partial_t\psi)^2/a_0^2$. For Sgr A*: $y \sim 10^{-4}$, shadow deviation $\sim 10^{-4}$ --- below ngEHT sensitivity.
\end{itemize}

\textbf{Grade: B (no change). GW constraints satisfied; shadow tests beyond current reach.}

%=============================================================================
\section{Agent 13: Particle Physics --- $\alpha$, Neutrinos}
%=============================================================================

\subsection{Mandate}
Track KATRIN final results, JUNO first results, $\alpha$ precision measurements.

\subsection{New Literature}

\begin{enumerate}

\item \textbf{KATRIN Collaboration (2025), Science adq9592} --- Direct neutrino mass: $m_\nu < 0.45$ eV (90\% CL) from 259 days / 36 million electrons. Factor 2 improvement over previous limit. Final dataset (1000 days by end 2025) targeting 0.3 eV. TRISTAN detector (keV sterile neutrinos) installation 2026. \textbf{Grade: A.}

DFD prediction: $m_\nu$ has no direct prediction, but the scalar field $\psi$ could affect neutrino mass via the disformal coupling. If neutrinos couple disformally: $m_\nu^\mathrm{eff} = m_\nu\,\sqrt{\Sigma(y)} = m_\nu/\sqrt{1+y}$. In the Newtonian regime ($y \ll 1$): negligible effect. Not yet testable.

\item \textbf{JUNO Collaboration (2025), arXiv:2511.14593} --- ``First measurement of reactor neutrino oscillations at JUNO.'' 59 days of data (Aug--Nov 2025). $\theta_{12}$ and $\Delta m^2_{21}$ measured with $1.6\times$ better precision than all previous experiments combined. Solar neutrino tension confirmed at previous $\sim$1.5$\sigma$ level. Mass hierarchy: requires $\sim$6 years for 3$\sigma$ determination. \textbf{Grade: A.}

\item \textbf{M\"uller group, Berkeley (2025), arXiv:2506.18328} --- ``The fine structure constant: a review of measurement results.'' Current best: $\alpha^{-1} = 137.035\,999\,206(11)$ from Rb recoil (81 ppt). Next-generation targeting 10 ppt via improved wavefront systematics. Paris Cs result consistent. \textbf{Grade: B.}

DFD prediction: $\alpha$ varies with $\psi$ as $\delta\alpha/\alpha = \kappa\,\psi$. With Th-229 nuclear clock $K = 5900$, the nuclear clock route is $\sim$5900$\times$ more sensitive than direct $\alpha$ measurement for DFD-type variations. Atom-interferometric $\alpha$ improvements are welcome but not the primary DFD test.

\end{enumerate}

\subsection{Neutrino Sector Summary}
\begin{itemize}
\item KATRIN: $m_\nu < 0.45$ eV. DFD has no strong neutrino prediction. \textbf{Neutral.}
\item JUNO: Solar tension confirmed. If resolved by new physics, could constrain scalar-neutrino coupling. \textbf{Watch.}
\item $\alpha$ precision: 81 ppt. Nuclear clock route ($K \sim 5900$) is the DFD-optimal channel.
\end{itemize}

\textbf{Grade: B.}

%=============================================================================
\section{Agent 14: Astrophysics --- Rotation Curves, EFE}
%=============================================================================

\subsection{Mandate}
Track rotation curve measurements, MOND tests in dwarf galaxies, external field effect.

\subsection{New Literature}

\begin{enumerate}

\item \textbf{Jiao et al.\ (2023), A\&A 678, A208; Ou et al.\ (2024)} --- Gaia DR3 Milky Way rotation curve: Keplerian decline starting at $\sim$16--19 kpc, ending at 26.5 kpc. Flat rotation curve rejected at 3$\sigma$. Total MW mass $M = 2.06\times10^{11}\,M_\odot$ --- order of magnitude less than standard DM halo. Baryonic fraction $\sim$1/3. \textbf{Grade: A.}

DFD/MOND interpretation: The declining rotation curve at large $R$ is \emph{expected} in MOND for the Milky Way because the enclosed baryonic mass is $\sim 6\times10^{10}\,M_\odot$ and the MOND radius $r_M = \sqrt{GM/a_0} \sim 30$ kpc. Beyond $r_M$, the rotation curve transitions from flat to declining. This is a \textbf{MOND success story}, not a problem.

\item \textbf{Mistele (2025), IJMPD 34, 2450034; arXiv:2409.17371} --- ``Implications of the Milky Way Declining Rotation Curve.'' MOG and Newtonian fits to Gaia DR3. MOG: $M \sim 1.3\times10^{11}\,M_\odot$; Newtonian: $M \sim 2\times10^{11}\,M_\odot$. \textbf{Grade: B.}

\item \textbf{Banik et al.\ (2024), A\&A 683, A148} --- ``Testing MOND using the dynamics of nearby stellar streams.'' Stream morphology in the Milky Way halo tests the gravitational potential. External field effect from the LMC and Sgr dwarf affects stream predictions. Mixed results: some streams fit MOND, others show tension. \textbf{Grade: B.}

\item \textbf{Frontiers (2025), Front.\ Astron.\ Space Sci.\ 1680387} --- ``A new empirical fit to galaxy rotation curves.'' New fitting function accommodating both flat and declining outer rotation curves. Tested on 175 galaxies from SPARC. \textbf{Grade: C.}

\item \textbf{Controversy (2025)} --- Recent critiques argue Jiao/Ou used disk-appropriate equations for spherical systems, and correct spherical modeling yields a flat rotation curve. If confirmed, this would remove the ``declining RC'' result. \textbf{Grade: B.} Important caveat: the declining MW RC may be a systematic artifact.

\end{enumerate}

\subsection{DFD Prediction Assessment}
DFD (via MOND limit) predicts:
\begin{itemize}
\item Flat rotation curves for $r < r_M = \sqrt{GM_b/a_0}$: \textbf{confirmed} for most galaxies.
\item Declining curves for $r > r_M$: \textbf{potentially confirmed} for MW (Jiao/Ou), but disputed.
\item EFE: dwarfs near massive hosts show reduced velocity dispersion: \textbf{confirmed} for several dwarfs.
\end{itemize}

\textbf{Grade: B+ (upgraded from B). MW declining RC, if confirmed, is a DFD/MOND prediction.}

%=============================================================================
\section{Agent 15: New Predictions from Disformal Upgrade}
%=============================================================================

\subsection{Mandate}
Identify novel tests arising from the disformal upgrade. Search for new observables.

\subsection{New Literature}

\begin{enumerate}

\item \textbf{Langlois (2025), arXiv:2412.04135} --- DHOST exact solutions include disformal Kerr BHs with measurable deviations from Kerr. \textbf{Grade: B.}

\item \textbf{Takahashi+ (2025), arXiv:2509.05027} --- GLPV $f^5$ induced GW spectrum. \textbf{Grade: B.}

\item \textbf{Barausse+ (2025), Nat.\ Astron.} --- ngEHT shadow measurements can test second post-Newtonian deviations. \textbf{Grade: B.}

\end{enumerate}

\subsection{Complete List of New/Updated Predictions (i26)}

\begin{center}
\begin{longtable}{clccc}
\toprule
\textbf{\#} & \textbf{Prediction} & \textbf{Channel} & \textbf{Grade} & \textbf{Status} \\
\midrule
\endhead
P1--P27 & (See i25 compilation) & Various & A/B/C & Various \\
P28 & $w_0 = -0.70\pm0.12$, $w_a = -0.85\pm0.25$ & Cosmo & A & Updated \\
P29 & Th-229 clock: $K = 5900\pm2300$ & Clock & A & On track \\
P30 & Velocity-dependent screening in mergers & Cluster & B & NEW \\
P31 & Gravitational slip $\eta < 1$ at $z < 0.5$ & Euclid & B & NEW \\
\bottomrule
\end{longtable}
\end{center}

\textbf{P30 (NEW --- Velocity-Dependent Screening):}
In merging galaxy clusters with high $\dot{\psi}$ (e.g., Bullet Cluster, El Gordo), the disformal factor $\Sigma(y) = 1/\sqrt{1+y}$ with $y = (\partial_t\psi)^2/a_0^2$ introduces a velocity-dependent correction to the effective gravitational coupling:
\begin{equation}
G_\mathrm{eff}(v) = G\,\left[1 + \mu(x)\right]\,\Sigma(y)
\end{equation}
For the Bullet Cluster collision velocity $v \sim 4700$ km/s and $\Phi_N \sim 10^{-5}c^2$: $\dot{\psi} \sim v\,\nabla\psi/c \sim v\,g_N/(c\,c) \sim 10^{-13}$ s$^{-1}$. Then $y = \dot{\psi}^2/a_0^2 \sim 10^{-7}$ --- too small to detect. \textbf{However}, for extreme mergers at $z > 1$ with $v > 0.01c$: $y$ could approach $10^{-3}$, giving a $\sim$0.05\% effect on the mass-to-light ratio. \textbf{Grade: B (long-term).}

\textbf{P31 (NEW --- Gravitational Slip):}
\begin{equation}
\eta_\mathrm{DFD} = \frac{\Phi}{\Psi} \approx 1 - \frac{y}{2(1+y)} \approx 1 - \frac{(\partial_t\psi)^2}{2\,a_0^2}
\end{equation}
On cosmological scales at $z < 0.5$ where dark energy effects are strongest, $y$ is determined by the Hubble flow: $\dot{\psi}_\mathrm{bg} \sim H_0\,\psi_\mathrm{bg}$. If $\psi_\mathrm{bg} \sim 10^{-5}$: $y \sim (H_0\,10^{-5})^2/a_0^2 \sim 10^{-12}$ --- undetectable.

\textbf{Reassessment}: P31 is only detectable if the scalar field background has significant time variation from dark energy dynamics ($\dot{\psi}_\mathrm{DE} \gg H_0\,\psi$). With DFD's dark energy EoS $w \sim -0.7$: $\dot{\psi}_\mathrm{DE} \sim \sqrt{3(1+w)}\,H_0\,M_{\rm Pl}/c^2$. This gives $y_\mathrm{DE} \sim 3(1+w) \sim 0.9$, so $\eta \sim 0.75$. \textbf{This is detectable by Euclid!}

\textbf{P31 Grade: A (if dark energy is DFD scalar). Euclid gravitational slip measurement is a critical test.}

%=============================================================================
\section{New Literature Master Table (Agents 1, 2, 3, 5, 13, 14, 15)}
%=============================================================================

\begin{longtable}{p{3cm}p{3cm}p{5cm}cc}
\toprule
\textbf{Reference} & \textbf{arXiv/DOI} & \textbf{Key Result} & \textbf{Agent} & \textbf{Grade} \\
\midrule
\endhead
Langlois 2025 & 2412.04135 & DHOST disformal gravity review & 1, 15 & A \\
Takahashi+ 2026 & 2601.09164 & 3rd-order ghost-free scalar-tensor & 1 & B \\
Kozak 2026 & 2603.08401 & Palatini disformal Horndeski & 1 & B \\
Takahashi+ 2025 & 2509.05027 & GLPV $f^5$ GW signature & 1, 5, 15 & B \\
MAGIS-100 2025 & Status report & 100m atom interferometer at Fermilab & 2 & B \\
Eot-Wash 2026 & 2602.02815 & Torsion balance EP test (Sun) & 2 & B \\
Qin+ 2024 & PDU 44, 101481 & Scalar field torsion balance limits & 2 & B \\
Biercuk+ 2025 & 2512.06333 & Quantum WEP test concept & 2 & C \\
Tiedau+ 2026 & Nature, 2507.01180 & Th-229 reproducibility $10^{-13}$ & 3 & A \\
Banerjee+ 2025 & Nat.\ Commun., 2407.17300 & $K = 5900\pm2300$ for Th-229 & 3 & A \\
JILA/Ye 2025 & Frequency measurement & 2,020,407,384,335(2) kHz & 3 & A \\
Thirolf+ 2025 & NSR nwaf083 & Th-229 developmental perspective & 3 & B \\
Ponce+ 2025 & PRL 134, 113801 & Th-229 temperature sensitivity & 3 & B \\
de Rham+ 2025 & 2511.19762 & GW speed constrains disformal & 5 & B \\
Barausse+ 2025 & Nat.\ Astron. & ngEHT shadow tests of gravity & 5, 15 & B \\
Saha+ 2025 & 2504.04331 & BH shadow in Lorentz gauge theory & 5 & C \\
Bernardo+ 2025 & GRG & BH labs review & 5 & C \\
KATRIN 2025 & Science & $m_\nu < 0.45$ eV, TRISTAN 2026 & 13 & A \\
JUNO 2025 & 2511.14593 & First reactor oscillation, 1.6$\times$ precision & 13 & A \\
M\"uller 2025 & 2506.18328 & $\alpha$ review, 81 ppt precision & 13 & B \\
Jiao+ 2023 & A\&A 678, A208 & MW declining RC from Gaia DR3 & 14 & A \\
Mistele 2025 & IJMPD 34, 2450034 & MW RC implications, MOG fits & 14 & B \\
Banik+ 2024 & A\&A 683, A148 & MOND stellar stream tests & 14 & B \\
\bottomrule
\end{longtable}

\end{document}
